\section{Implementation}

Our inference system is implemented on top of llama.cpp~\cite{llama.cpp} with approximately 7K lines of code in C/C++ and inline assembly. We use the LLVM toolchain in the Hexagon SDK (version 6.0.0.2) to generate code for Hexagon NPUs. We especially note that our system has no dependency on Qualcomm's QNN, avoiding inflexible static fixed-shape computation graphs.

Our implementation mainly consists of two modules: one module is the operator library for the Hexagon NPU, which is compiled into an independent Hexagon DSP shared object; the other module is integrated with llama.cpp on the CPU side. The NPU operator library implements computation kernels, power management, hardware resource management, and a computation thread pool. We add a Hexagon NPU backend to llama.cpp, leveraging \texttt{rpcmem} shared memory as the underlying buffer type. \texttt{rpcmem} is a wrapper for the kernel dmabuf memory and supports the sharing of physical memory between the CPU and the NPU. The related allocation, deallocation, and mapping interfaces are provided by \texttt{libcdsprpc.so} in the Android system’s vendor libraries. By utilizing shared memory buffers, we not only eliminate unnecessary inter-processor data copy but also reuse the existing memory management system as much as possible. In addition, we are able to schedule the operators that have not been implemented on the NPU to run on the CPU, achieving seamless integrations with upper-layer applications.

During the backend initialization phase, we call the FastRPC~\cite{fastrpc} facility of the Hexagon SDK to start the remote NPU session and initialize an area of shared memory for communication. On the NPU side, a thread continuously polls in this shared-memory area to receive computation requests from the CPU. Compared to the default RPC implementation, communication through shared memory can have a lower latency. We note that after the CPU writes data to the shared memory, the NPU will not automatically invalidate the cache of the corresponding area as there is only one-way coherence between the CPU and the NPU on the Snapdragon SoC. Therefore, we manually clear the cache before NPU polls. Similar cache maintenance operations are also required for shared buffers containing model activations.


% \subsection{De-quantization based mixed-precision GEMM on NPU}

% Although Qualcomm Hexagon NPU supports native mixed INT4/INT16 GEMM (also referred to as W4A16, where weight is quantized to 4 bits and activation is quantized to 16 bits integer values), we choose to use dequantization-based methods and utilize FP16 GEMM components due to the following reasons:

% 1. Hexagon NPU's native W4A16 GEMM suffers from accuracy loss. The W4A16 GEMM hardware component is designed to use the per-channel quantization scheme, i.e. weight elements in a row/column share one floating-point scale and zero-point. However, the row/column size of attention projection and FFN weight matrices in transformers model can easily reach a few thousands, forming large quantization groups and resulting in significant quantization errors. For comparison, the widely adopted Q4\_0 and Q4\_K 4-bit quantization schemes used in llama.cpp have a much finer-grained quantization group of 32 elements. Some works (e.g. powerinfer-2) have observed the accuracy issue of Qualcomm's official implementation QNN (Qualcomm AI Engine Direct) due to the quantization scheme. \sout{To fit in Hexagon NPU's hardware design, advanced quantization methods like AWQ and QuaRot(?, cite) need to be applied to decrease quantization errors, some of which require extra calibration process, reducing the usability of NPU's native mixed-precision GEMM.} \hzx{AWQ is a weight-only quantization method and its calibration process is much simpler compared to other methods. Don't mix weight-only quantization and fully quantized scheme (both weight and activation are quantized). Maybe extra efforts are needed to integrate fully-quantized schemes into existing pipelines?}

% 2. Most of the Hexagon NPU's high-throughput GEMM instructions are not publicly available. 

% \subsubsection{Weight dequantization using HVX}

% Since Hexagon NPU's HMX unit is dedicated to executing matrix multiplication, we use NPU's vector SIMD components (Hexagon Vector eXtension, HVX for short) to perform weight dequantization. Unfortunately, HVX has very limited computation capabilities compared to HMX, \sout{for example, only about 50 GFLOPS. What makes things worse is that HVX mainly features integer SIMD operations, only supporting basic floating-point operations like adding and multiplying} \sout{(floating-point vector division is not supported)}. Moreover, the HVX vector registers have 1024 bit width, each can hold 128 bytes or 64 FP16 elements. This leads to computation waste or performance drop when operating on fine-grained quantization groups: we either compute with vector registers containing padded elements or use extra instruction to combine operands into one vector register.

% To overcome these hardware constraints, we reorder quantization groups to exploit the width of the HVX vector register and use table lookup instructions instead of costly integer-float conversion. Taking the Q4\_0 quantization scheme as an example, 32 elements in one quantization group share the same FP16 scale, and the quantized values take up 16 bytes. The scales and quantized values are interleaved in memory. We repack 8 quantization groups so that 8 * 16 quantized values become continuous and fit in a single HVX register. The vlut16 instruction is capable of performing a table lookup in a table of 16 elements for each 8-bit index in a source register. Using vlut16 instructions, we directly transform 4-bits quantized values into [-8, 7) FP16 values (\hzx{Now we also support IQ4\_NL non-linear quantization scheme where INT4 values are mapped to values with uneven spaces}), as well as duplicate the packed scales to correct positions for further multiplication without incurring additional permutations. 

% We note that our Q4\_0 and Q8\_0 weight dequantization procedure is bounded by HVX compute, not memory bandwidth. This characteristic is quite unique when compared to modern GPUs. Our initial attempt to use popular Q4\_K quantization scheme failed with very poor performance as element group scales are quantized into 6-bits integers, unpacking and converting them into floating-point elements using HVX introduces significant overhead. 

% \subsubsection{Overlapping weight transfer and computation with asynchronous DMA}

% Hexagon NPU's memory system also exhibits special characteristics: the HVX unit accepts memory operands from either L2 cache or VTCM (Vector Tightly Coupled Memory). The HVX unit is designed to bypass the L1 cache to avoid cache pollution of scalar cores. VTCM serves as a software-managed L1 cache and has a typical size of 8 MiB. Using VTCM is preferred because it provides low-latency access and consumes less energy.

% Our experimental results show that using regular HVX vector memory accesses to transfer data from DDR memory to VTCM (datapath: DDR -> L2 cache -> VTCM) with multiple vector cores and software-based asynchronous L2 cache prefetch cannot fully utilize the DDR memory bandwidth, only reaching 26 GB/s on a 8gen3 device.

% Hexagon NPUs also provide asynchronous DMA functionality to enable fast movement between DDR memory and VTCM, achieving near 60 GB/s bandwidth. We leverage this feature to asynchronously prefetch weight chunks into VTCM, completely hiding data loading latency.

% \subsection{Fast Inter-Processor Communication via Shared Memory}

% Since HMX only supports matrix multiplication/convolution and the general-purpose floating-point computing capability of HVX is relatively weak, not all parts of the transformer models are suitable for computation on Hexagon NPU. For example, HVX hardware does not support native exponential calculation; therefore, it is usually necessary to expand the exponential to a polynomial using the Taylor series for calculation on Hexagon NPU. \sout{However, the floating-point multiply/add operations of HVX have only about 1/4 of the throughput compared to integer operations.} This makes the computation of the softmax in Attention and the SiLU activation function on the NPU incur a high overhead. In addition, the limited number of floating-point instructions of HVX and the extremely wide vector registers also increase the difficulty and cost of optimizing the NPU operators.

% A reasonable approach is to offload suitable operators such as matrix multiplication to the NPU and compute the remaining operators on the CPU, so as to achieve a high degree of flexibility and a reasonable performance trade-off. However, such a scheme must take into account the cost of inter-processor communication between the CPU and the NPU. The default communication method between the Hexagon NPU and the CPU is FastRPC. In general, the processor core will fall asleep while waiting for the response from the peer and will rely on cross-processor interrupts to wake up. Our measurement shows that once the processor falls asleep, it will lead to high communication latency. The latency of communication via RPC is often above 100 microseconds, which may be several times the actual execution time of the operators. The high communication latency of the default RPC poses difficulties for flexible task partitioning between the CPU and the NPU.

% In fact, considering the characteristic that the mobile SoC shares the global DDR memory, we doubt the rationality of RPC as the inter-processor communication method. The overheads related to the RPC kernel path, serialization/deserialization may be unnecessary in a shared memory system. Therefore, we propose to achieve efficient communication between the CPU and the NPU through shared memory.
